%% This file is based on https://github.com/tudace/tuda_latex_templates by Marei Peischl <marei@peitex.de>
%%
% !TeX program = lualatex
%%

\documentclass[
	language=english, % Zu ngerman ändern bei deutschen Arbeiten
	thesis={type=msc}, % Type ändern auf ... sta: Studienarbeit, bsc: Bachelor, msc: Master
%% Die folgenden Parameter sollten nicht geändert werden
	class=report,% Basisdokumentenklasse. Wählt die Korrespondierende KOMA-Script Klasse
	ruledheaders=section,% Ebene bis zu der die Überschriften mit Linien abgetrennt werden, vgl. DEMO-TUDaPub
	accentcolor=3d,% Auswahl der Akzentfarbe
	marginpar=false,% Kopfzeile und Fußzeile erstrecken sich nicht über die Randnotizspalte
	parskip=half-,% Absatzkennzeichnung durch Abstand vgl. KOMA-Script
	fontsize=11pt,% Basisschriftgröße laut Corporate Design ist mit 9pt häufig zu klein
	IMRAD=false, % Zusätzliche Metadaten nach Wunsch der Universitätsbibliothek deaktivieren
%	pdfa=false % tudapub hat einen Workaround für pdfa compilation https://github.com/tudace/tuda_latex_templates/issues/321 Falls kein pdfa notwendig ist: pdfa=false nutzen und das \Metadata Command im Document head entfernen
]{tudapub}

%%%%%%%%%%%%%%%%%%%
% Sprachanpassung & Verbesserte Trennregeln
%%%%%%%%%%%%%%%%%%%
\usepackage[ngerman, main=english]{babel} % ngerman und englisch bei deutschen Arbeiten tauschen
\usepackage[english]{datetime2} % german und english bei der Arbeit anpassen
\usepackage[autostyle]{csquotes} % Anführungszeichen vereinfacht
\usepackage{microtype}

%%%%%%%%%%%%%%%%%%%
% Literaturverzeichnis
%%%%%%%%%%%%%%%%%%%
\usepackage[
    backend=biber,
    bibencoding=utf-8,
    language=auto,
    sorting=nyt,
    defernumbers=true,
    style=ieee,
    maxbibnames=100,
    mincitenames=1,
    maxcitenames=2,
  ]{biblatex} % Literaturverzeichnis
\bibliography{bibliography}

%%%%%%%%%%%%%%%%%%%
% Abkürzungen (Akronym Referenzen durch \ac{SG} oder \acp{SG} für Singular oder Plural Formen
%%%%%%%%%%%%%%%%%%%
\usepackage{acro}
\usepackage{mfirstuc}% provides \capitalisewords
\acsetup{list/template=longtable}
\input{acronyms}

%%%%%%%%%%%%%%%%%%%
% Paketvorschläge Tabellen
%%%%%%%%%%%%%%%%%%%
%\usepackage{array}     % Basispaket für Tabellenkonfiguration, wird von den folgenden automatisch geladen
\usepackage{tabularx}   % Tabellen, die sich automatisch der Breite anpassen
\usepackage{longtable} % Mehrseitige Tabellen
%\usepackage{xltabular} % Mehrseitige Tabellen mit anpassarer Breite
\usepackage{booktabs}   % Verbesserte Möglichkeiten für Tabellenlayout über horizontale Linien

%%%%%%%%%%%%%%%%%%%
% Paketvorschläge Mathematik
%%%%%%%%%%%%%%%%%%%
%\usepackage{mathtools} % erweiterte Fassung von amsmath
%\usepackage{amssymb}   % erweiterter Zeichensatz
%\usepackage{siunitx}   % Einheiten

%%%%%%%%%%%%%%%%%%%
% Weitere Paketvorschläge
%%%%%%%%%%%%%%%%%%%
\usepackage{pifont}% Zapf-Dingbats Symbole
\usepackage{subcaption}
\usepackage{framed}

\input{macros}

\begin{document}

\title{Title of Thesis}
\subtitle{Subtitle}
\author[F. Surname]{Forename Surname} % optionales Argument ist die Signatur
\reviewer{Dr.-Ing. (habil.) Stefan Göbel \and Supervisor} % Gutachter
\department{etit}
\institute{Serious Games}

\Metadata{
    author=Forename Surname,
    title=Title,
}

\submissiondate{\today}
\examdate{\today}

% Hinweis zur Lizenz:
% TUDa-CI verwendet momentan die Lizenz CC BY-NC-ND 2.0 DE als Voreinstellung.
% Die TU Darmstadt hat jedoch die Empfehlung von dieser auf die liberalere
% CC BY 4.0 geändert. Diese erlaubt eine Verwendung bearbeiteter Versionen und
% die kommerzielle Nutzung.
% TUDa-CI wird im nächsten größeren Release ebenfalls diese Anpassung vornehmen.
% Aus diesem Grund wird empfohlen die Lizenz manuell auszuwählen.
% \tuprints{urn=1234,printid=12345,doi=10.25534/tuprints-1234,license=cc-by-4.0}
% To see furhter information on the license option in English, remove the license= key and pay attention to the warning & help message.

\maketitle
\affidavit % oder \affidavit[digital] falls eine rein digitale Abgabe vorgesehen ist.
\tableofcontents
\input{abstract}

%%%%%%%%%%%%%%%%%%%
% Content
%%%%%%%%%%%%%%%%%%%
\input{content}

\clearpage
\chapter*{List of Acronyms}
\printacronyms[heading={none}]

\printbibliography

\appendix
\input{appendix}

\end{document}
